\scanpage{111}
est \unsure{claire|clair} sur l'expression 155, \deleted{\unsure{les}} $J'\subset J$, Considérons \deleted{le plus petit}\inserted{un} élément \deleted{$e_{J'}$}\inserted{$x$} de $J'$, \inserted{$\ne e_J$} \deleted{qui} est de codim $\delta$ \inserted{($\unsure{\le|\ne}d$)} et considérons \inserted{d'abord} le composé
\[
(158)\quad
\diagram{\begin{tikzcd}[ampersand replacement=\&, column sep=large,row sep=small]
D_J\arrow[r]\arrow[dr,draw=none,"{\deleted{\searrow}}" description]\&D_{J'}\arrow[r,"f_x"]\&D_\delta\\
\&\deleted{D_{d,\delta}}\arrow[ur,draw=none,"{\deleted{\nearrow\,\unsure{\mathrm{can.}}}}" description]\&
\end{tikzcd}}
\]
\margin{\inserted{$f_x$ défini par fonctorialité par l'inclusion $\{x\}\subset J'$}}
on va le mettre en relation avec l'hom.
\[
(159)\qquad D_J\longrightarrow D_{d,\delta}
\]
suivant (où $D_{d,\delta}=D_{\Delta(d,\delta)}$, $\Delta(d,\delta)$ étant le drapeau \deleted{type} de type $d,\delta$, i.e. l'ens. ordonné \deleted{avec} fonction codim
\diagram{\begin{tikzcd}[ampersand replacement=\&, column sep=small]\underset{d}{\bullet}\arrow[r,dash]\&\underset{\delta}{\bullet}\end{tikzcd}}
(si $\delta<d$!) \deleted{\unsure{resp.} \diagram{\begin{tikzcd}\underset{d=\delta}{\bullet}\end{tikzcd}}}
\deleted{(\unsure{si|et} $\delta=1$)}, on a, en appliquant (155) au cas de $\inserted{J=}\Delta(d,\delta)$, \deleted{au lieu de} \deleted{\unsure{l'}} l'isom canonique
\[
(160)\qquad D_{d,\delta}\simeq\unsure{\widetilde D'_d|\widetilde{\widetilde D}_d}\mathbin{\unsure{\wedge|\times}}_{(\unsure{\mathfrak S|\mathcal S}_d)_{D_d}}\mathcal P_\delta(I_{\inserted{d}}\deleted{(d)})
\]
D'autre part, on a (avec la notation (154) pour $H(J)$), \deleted{pour} l'élément $\unsure{x}\in J$ fixé de $J$, \inserted{\unsure{de codim} $d'=\operatorname{cod}(e')$,} une application
\[
(161)\quad\diagram{\begin{tikzcd}[ampersand replacement=\&, row sep=small]
H(J)\arrow[r]\&\mathcal P_\delta(I_d)\\
h\arrow[r,mapsto]\&h(\unsure{x})
\end{tikzcd}}
\]
\margin{110}
